% Part: first-order-logic
% Chapter: syntax
% Section: intro-syntax

\documentclass[../../../include/open-logic-section]{subfiles}

\begin{document}

\olfileid{fol}{syn}{itx}

\olsection{ভূমিকা}

প্রথম-ক্রম যুক্তিবিদ্যার তত্ত্ব ও অধিতত্ত্ব গড়ে তুলতে হলে প্রথমেই তার
রাশিগুলির সংকেতবিন্যাস ও অর্থতত্ত্ব সংজ্ঞায়িত করতে হবে। প্রথম-ক্রম
যুক্তিবিদ্যার রাশি হলো পদ ও !!{formula}s। !!{variable}s,
!!{constant}s ও !!{function}s থেকে পদ গঠিত হয়। আবার !!^{formula}s,
!!{predicate}s-এর
সঙ্গে পদ মিলিয়ে সবচেয়ে ছোট, “পরমাণু” !!{formula}s গঠিত হয়; তারপর
যৌক্তিক সংযোজক ও পরিমাণসূচক ব্যবহার করে পরমাণু !!{formula}s থেকে আরও
জটিল সূত্র গঠন করা যায়। গঠনের নিয়মগুলি নির্ধারণের নানা উপায় আছে;
আমরা তার মধ্যে কেবল একটি উপায় দিচ্ছি। অন্য ব্যবস্থায় ভিন্ন সংকেত
বেছে নেওয়া হতে পারে, ভিন্ন সংযোজকগুচ্ছকে মৌলিক ধরা হতে পারে,
বন্ধনী ভিন্নভাবে ব্যবহার করা হতে পারে (কিংবা তথাকথিত পোলিশ সংকেতরীতির
মতো একেবারেই ব্যবহার না-ও করা হতে পারে)। তবে সব পদ্ধতির সাধারণ বৈশিষ্ট্য
হলো, গঠনের নিয়মগুলি পদ ও !!{formula}s-এর সেটকে \emph{আরোহীভাবে}
সংজ্ঞায়িত করে। নিয়মগুলি যথাযথ হলে, প্রত্যেকটি রাশি মূলত একটিমাত্র
উপায়েই সেই নিয়ম মেনে গঠিত হতে পারে। আরোহী সংজ্ঞা থেকে রাশিগুলি
\emph{একক পাঠযোগ্য} হওয়ার ফলে একই পদ্ধতিতে—আরোহী সংজ্ঞার সাহায্যে—আমরা
সেগুলির অর্থও দিতে পারি।

\end{document}
